% GENERATED FILE. Edit canonical lessons, then run npm run generate:books.
% canonical-source-hash: fnv1a64:50673d02fa4622e3
% canonical-lessons: SA-C290-dasyuh, SA-C290-lunthati, SA-C290-akramati, SA-C290-abharanam, SA-C290-balat

\chapter{Robbers on the Road}
\label{ch:sa-robbers-on-the-road}
\hlchaptermodality{290}

\begin{chapteropening}
\textbf{By the end of this chapter:} \emph{I can tell what robbers did and say what happens if they come.}

Complete the last lesson of chapter 290: I can tell what robbers did and say what happens if they come.
\end{chapteropening}

\section[\texorpdfstring{\sk{दस्युः} (dasyuḥ)}{दस्युः (dasyuḥ)}]{\sk{दस्युः} (dasyuḥ) — a robber; a bandit}

\label{lesson:SA-C290-dasyuh}



\begin{quote}
Before the new one: say the Sanskrit for a trace, then the Sanskrit for before one's eyes.
\end{quote}

\subsection*{You'll want to know: \sk{दस्युः}}
\textbf{\sk{दस्युः}} — \emph{dasyuḥ} — \textquotedblleft{}a robber; a bandit\textquotedblright{}.

\begin{quote}
\sk{वने} \sk{दस्युः} \sk{वसति।}
\end{quote}

\begin{quote}
\emph{vane dasyuḥ vasati.}
\end{quote}

\begin{quote}
A robber lives in the forest.
\end{quote}

A masculine noun.

\begin{cousinweb}
\textbf{\sk{दस्युः}} (\emph{dasyuḥ}) comes from the root √das, \textquotedblleft{}to lay waste\textquotedblright{}; MW's older sense is an enemy of the gods, and it also gives \emph{dasyu} as a robber (MW).
\end{cousinweb}

\begin{grammarlens}[title={What you've built}]
One more word for telling what happened.
\end{grammarlens}

\subsection*{Your turn}
\begin{itemize}
\raggedright
  \item \emph{Say it:} \emph{dasyuḥ}
  \item \emph{Say it:} \emph{dasyuḥ}, once more
  \item \emph{Say it:} \emph{vane dasyuḥ vasati.} — A robber lives in the forest.
  \item \emph{Recall:} say the Sanskrit for a trace, then the Sanskrit for before one's eyes, then say \emph{dasyuḥ} again
\end{itemize}

\begin{tcolorbox}[breakable,colback=teal!4,colframe=teal!35!black,title={Before you move on}]
What does \sk{दस्युः} mean? (\textquotedblleft{}A robber; a bandit\textquotedblright{}.) Say it once more.
\end{tcolorbox}
\section[\texorpdfstring{\sk{लुण्ठति} (luṇṭhati)}{लुण्ठति (luṇṭhati)}]{\sk{लुण्ठति} (luṇṭhati) — plunders}

\label{lesson:SA-C290-lunthati}



\begin{quote}
Before the new one: say the Sanskrit for before one's eyes, then the Sanskrit for a robber.
\end{quote}

\subsection*{You'll want to know: \sk{लुण्ठति}}
\textbf{\sk{लुण्ठति}} — \emph{luṇṭhati} — \textquotedblleft{}plunders\textquotedblright{}.

\begin{quote}
\sk{यदि} \sk{दस्युः} \sk{आगच्छति}, \sk{तर्हि} \sk{ग्रामं} \sk{लुण्ठति।}
\end{quote}

\begin{quote}
\emph{yadi dasyuḥ āgacchati, tarhi grāmaṃ luṇṭhati.}
\end{quote}

\begin{quote}
If a robber comes, then he plunders the village.
\end{quote}

Past (he …): \sk{सः} \sk{ग्रामं} \sk{लुण्ठितवान्}, \emph{saḥ grāmaṃ luṇṭhitavān}.

\begin{cousinweb}
\textbf{\sk{लुण्ठति}} (\emph{luṇṭhati}) comes from the root √luṇṭh, which MW gives with several senses, among them to rob or plunder (MW).
\end{cousinweb}

\begin{grammarlens}[title={What you've built}]
One more word for telling what happened.
\end{grammarlens}

\subsection*{Your turn}
\begin{itemize}
\raggedright
  \item \emph{Say it:} \emph{luṇṭhati}
  \item \emph{Say it:} \emph{luṇṭhati}, once more
  \item \emph{Say it:} \emph{yadi dasyuḥ āgacchati, tarhi grāmaṃ luṇṭhati.} — If a robber comes, then he plunders the village.
  \item \emph{Recall:} say the Sanskrit for before one's eyes, then the Sanskrit for a robber, then say \emph{luṇṭhati} again
\end{itemize}

\begin{tcolorbox}[breakable,colback=teal!4,colframe=teal!35!black,title={Before you move on}]
What does \sk{लुण्ठति} mean? (\textquotedblleft{}Plunders\textquotedblright{}.) Say it once more.
\end{tcolorbox}
\section[\texorpdfstring{\sk{आक्रमति} (ākramati)}{आक्रमति (ākramati)}]{\sk{आक्रमति} (ākramati) — attacks}

\label{lesson:SA-C290-akramati}



\begin{quote}
Before the new one: say the Sanskrit for a robber, then the Sanskrit for plunders.
\end{quote}

\subsection*{You'll want to know: \sk{आक्रमति}}
\textbf{\sk{आक्रमति}} — \emph{ākramati} — \textquotedblleft{}attacks\textquotedblright{}.

\begin{quote}
\sk{दस्युः} \sk{पथिकम्} \sk{आक्रमति।}
\end{quote}

\begin{quote}
\emph{dasyuḥ pathikam ākramati.}
\end{quote}

\begin{quote}
The robber attacks the traveller.
\end{quote}

Past (he …): \sk{सः} \sk{पथिकम्} \sk{आक्रान्तवान्}, \emph{saḥ pathikam ākrāntavān}.

\begin{cousinweb}
\textbf{\sk{आक्रमति}} (\emph{ākramati}) is built from \emph{ā}, \textquotedblleft{}towards\textquotedblright{}, and the root √kram, \textquotedblleft{}to step, stride\textquotedblright{}; MW gives \emph{ā-kram} as to step near, seize and attack (MW).
\end{cousinweb}

\begin{grammarlens}[title={What you've built}]
One more word for telling what happened.
\end{grammarlens}

\subsection*{Your turn}
\begin{itemize}
\raggedright
  \item \emph{Say it:} \emph{ākramati}
  \item \emph{Say it:} \emph{ākramati}, once more
  \item \emph{Say it:} \emph{dasyuḥ pathikam ākramati.} — The robber attacks the traveller.
  \item \emph{Recall:} say the Sanskrit for a robber, then the Sanskrit for plunders, then say \emph{ākramati} again
\end{itemize}

\begin{tcolorbox}[breakable,colback=teal!4,colframe=teal!35!black,title={Before you move on}]
What does \sk{आक्रमति} mean? (\textquotedblleft{}Attacks\textquotedblright{}.) Say it once more.
\end{tcolorbox}
\section[\texorpdfstring{\sk{आभरणम्} (ābharaṇam)}{आभरणम् (ābharaṇam)}]{\sk{आभरणम्} (ābharaṇam) — a piece of jewellery}

\label{lesson:SA-C290-abharanam}



\begin{quote}
Before the new one: say the Sanskrit for plunders, then the Sanskrit for attacks.
\end{quote}

\subsection*{You'll want to know: \sk{आभरणम्}}
\textbf{\sk{आभरणम्}} — \emph{ābharaṇam} — \textquotedblleft{}a piece of jewellery\textquotedblright{}.

\begin{quote}
\sk{दस्युः} \sk{आभरणं} \sk{हरति।}
\end{quote}

\begin{quote}
\emph{dasyuḥ ābharaṇaṃ harati.}
\end{quote}

\begin{quote}
The robber takes the jewellery.
\end{quote}

A neuter noun.

\begin{cousinweb}
\textbf{\sk{आभरणम्}} (\emph{ābharaṇam}) is built from \emph{ā} and the root √bhṛ, \textquotedblleft{}to bear, carry\textquotedblright{}; MW gives \emph{ābharaṇa} as decorating and as an ornament such as jewels (MW).
\end{cousinweb}

\begin{grammarlens}[title={What you've built}]
One more word for telling what happened.
\end{grammarlens}

\subsection*{Your turn}
\begin{itemize}
\raggedright
  \item \emph{Say it:} \emph{ābharaṇam}
  \item \emph{Say it:} \emph{ābharaṇam}, once more
  \item \emph{Say it:} \emph{dasyuḥ ābharaṇaṃ harati.} — The robber takes the jewellery.
  \item \emph{Recall:} say the Sanskrit for plunders, then the Sanskrit for attacks, then say \emph{ābharaṇam} again
\end{itemize}

\begin{tcolorbox}[breakable,colback=teal!4,colframe=teal!35!black,title={Before you move on}]
What does \sk{आभरणम्} mean? (\textquotedblleft{}A piece of jewellery\textquotedblright{}.) Say it once more.
\end{tcolorbox}
\section[\texorpdfstring{\sk{बलात्} (balāt)}{बलात् (balāt)}]{\sk{बलात्} (balāt) — by force}

\label{lesson:SA-C290-balat}



\begin{quote}
Before the new one: say the Sanskrit for attacks, then the Sanskrit for a piece of jewellery.
\end{quote}

\subsection*{You'll want to know: \sk{बलात्}}
\textbf{\sk{बलात्}} — \emph{balāt} — \textquotedblleft{}by force\textquotedblright{}.

\begin{quote}
\sk{दस्युः} \sk{बलात्} \sk{धनं} \sk{गृह्णाति।}
\end{quote}

\begin{quote}
\emph{dasyuḥ balāt dhanaṃ gṛhṇāti.}
\end{quote}

\begin{quote}
The robber takes the money by force.
\end{quote}

\begin{cousinweb}
\textbf{\sk{बलात्}} (\emph{balāt}) is built on \emph{bala}, \textquotedblleft{}strength, force\textquotedblright{}: MW records its ablative \emph{balāt} as meaning forcibly, against one's will (MW).
\end{cousinweb}

\begin{grammarlens}[title={What you've built}]
One more word for telling what happened.
\end{grammarlens}

\subsection*{Your turn}
\begin{itemize}
\raggedright
  \item \emph{Say it:} \emph{balāt}
  \item \emph{Say it:} \emph{balāt}, once more
  \item \emph{Say it:} \emph{dasyuḥ balāt dhanaṃ gṛhṇāti.} — The robber takes the money by force.
  \item \emph{Recall:} say the Sanskrit for attacks, then the Sanskrit for a piece of jewellery, then say \emph{balāt} again
\end{itemize}

\begin{tcolorbox}[breakable,colback=teal!4,colframe=teal!35!black,title={Before you move on}]
What does \sk{बलात्} mean? (\textquotedblleft{}By force\textquotedblright{}.) Say it once more.
\end{tcolorbox}
